\section{Conceptos fundamentales del aprendizaje por refuerzo}

Esta sección presenta en detalle los cinco conceptos fundamentales del aprendizaje por refuerzo: Agent (agente), Environment (entorno), Action (acción), Observation/State (observación/estado), Reward (recompensa), así como la relación de interacción entre ellos.

\subsection{Agent y Environment}

\begin{figure}[H]
\centering
\includegraphics[width=0.85\textwidth]{slides-images/slide-11.jpg}
\caption{El circuito de interacción central de RL: el Agent interactúa con el Environment mediante Action y recibe Observation\protect\footnotemark}
\end{figure}
\footnotetext{Diapositiva 11.}

\begin{itemize}
    \item \textbf{Agent} (agente): el sujeto de decisión dentro del sistema, encargado de tomar acciones con base en la observación actual. Por ejemplo: el personaje de Mario en Super Mario, un dron, el controlador de un vehículo autónomo
    \item \textbf{Environment} (entorno): el mundo en el que existe y opera el Agent. Las reglas del entorno definen cómo las acciones del Agent afectan el estado del mundo
\end{itemize}

\subsection{Action y Action Space}

\textbf{Action} (acción) es la operación que el Agent puede ejecutar en el entorno, denotada como $a_t$ (la acción en el instante $t$). El conjunto de todas las acciones posibles se denomina \textbf{Action Space} (espacio de acciones) $\mathcal{A}$.

El espacio de acciones puede ser:
\begin{itemize}
    \item \textbf{Discreto} (Discrete): por ejemplo, las cuatro direcciones arriba, abajo, izquierda y derecha
    \item \textbf{Continuo} (Continuous): por ejemplo, el ángulo del volante (cualquier valor real)
\end{itemize}

\subsection{Observation y State}

Después de que el Agent toma una acción, el Environment devuelve una nueva \textbf{observación} (Observation), que refleja el cambio producido en el entorno por dicha acción. La situación actual percibida por el Agent se denomina \textbf{State} (estado), denotado como $s_t$.

\begin{figure}[H]
\centering
\includegraphics[width=0.85\textwidth]{slides-images/slide-12.jpg}
\caption{Cambio de estado: después de que el Agent actúa, el entorno devuelve el nuevo estado $s_{t+1}$\protect\footnotemark}
\end{figure}
\footnotetext{Diapositiva 12.}

\subsection{Reward y Return}

\textbf{Reward} (recompensa) $r_t$ es la señal escalar que el entorno devuelve al Agent en cada paso de tiempo, y que indica qué tan «bueno» es el estado actual. El objetivo del Agent es maximizar el \textbf{retorno total} (Total Reward / Return):

$$R_t = \sum_{i=t}^{\infty} r_i$$

\begin{itemize}
    \item $R_t$: el retorno total a partir del instante $t$
    \item $r_i$: la recompensa inmediata obtenida en el instante $i$
\end{itemize}

\begin{figure}[H]
\centering
\includegraphics[width=0.85\textwidth]{slides-images/slide-14.jpg}
\caption{El circuito de interacción completo de RL: el Agent envía Action, el Environment devuelve State y Reward\protect\footnotemark}
\end{figure}
\footnotetext{Diapositiva 14.}

\subsection{Retorno descontado (Discounted Return)}

En las aplicaciones reales, no todas las recompensas tienen la misma importancia: las \textbf{recompensas cercanas} suelen valer más que las lejanas (porque el futuro conlleva incertidumbre). Por ello se introduce un \textbf{factor de descuento} $\gamma \in [0, 1)$ para atenuar las recompensas futuras:

$$R_t = \sum_{i=t}^{\infty} \gamma^{i-t} r_i = r_t + \gamma r_{t+1} + \gamma^2 r_{t+2} + \cdots$$

\begin{itemize}
    \item $\gamma$: el factor de descuento (Discount Factor), que controla cuánto peso se le da a las recompensas futuras
    \item Cuanto más cercano a 0 es $\gamma$, más «corto de vista» es el Agent, y solo atiende a la recompensa inmediata
    \item Cuanto más cercano a 1 es $\gamma$, más «largo de vista» es el Agent, y presta más atención al retorno de largo plazo
\end{itemize}

\begin{figure}[H]
\centering
\includegraphics[width=0.85\textwidth]{slides-images/slide-16.jpg}
\caption{Retorno descontado: el factor de descuento $\gamma$ atenúa el peso de las recompensas futuras\protect\footnotemark}
\end{figure}
\footnotetext{Diapositiva 16.}

\begin{warningbox}{¿Por qué se necesita un factor de descuento?}
Si no se usa un factor de descuento ($\gamma = 1$), en un horizonte de tiempo infinito el retorno total puede tender a infinito, lo que vuelve irresoluble el problema de optimización. Además, el factor de descuento también refleja la intuición del mundo real de que «más vale pájaro en mano que ciento volando»: las recompensas futuras conllevan incertidumbre, mientras que la recompensa presente y segura vale más.
\end{warningbox}

\subsection{Q-function (función de valor estado-acción)}

La \textbf{Q-function} es uno de los conceptos más centrales del aprendizaje por refuerzo. Mide el \textbf{retorno total futuro esperado} que el Agent puede obtener al ejecutar la acción $a$ en el estado $s$:

$$Q(s_t, a_t) = \mathbb{E}[R_t \mid s_t, a_t]$$

\begin{figure}[H]
\centering
\includegraphics[width=0.85\textwidth]{slides-images/slide-18.jpg}
\caption{Definición de la Q-function: el retorno total futuro esperado\protect\footnotemark}
\end{figure}
\footnotetext{Diapositiva 18.}

\begin{importantbox}{El significado central de la Q-function}
La pregunta que responde $Q(s, a)$ es: «si tomo la acción $a$ en el estado $s$ y, a partir de ahí, sigo siempre la política óptima, ¿cuánto retorno total puedo obtener?». No solo considera la recompensa inmediata de la acción actual, sino también su efecto sobre todos los \textbf{estados y recompensas futuros}.
\end{importantbox}

\subsection{Policy (política)}

Con la Q-function, el Agent puede establecer una \textbf{política} (Policy), denotada como $\pi(s)$. La política óptima $\pi^*(s)$ consiste en elegir, en cada estado, la acción que maximiza la Q-function:

$$\pi^*(s) = \arg\max_a Q(s, a)$$

\begin{figure}[H]
\centering
\includegraphics[width=0.85\textwidth]{slides-images/slide-19.jpg}
\caption{Cómo usar la Q-function para establecer una política: elegir la acción que maximiza $Q(s,a)$\protect\footnotemark}
\end{figure}
\footnotetext{Diapositiva 19.}

\subsection{Resumen de la sección}

El marco central del aprendizaje por refuerzo puede resumirse como un proceso de interacción en circuito cerrado: el Agent observa el estado actual $s_t$, elige la acción $a_t$ según la política $\pi$, el entorno devuelve el nuevo estado $s_{t+1}$ y la recompensa $r_t$, y el objetivo del Agent es encontrar la política óptima $\pi^*$ que maximice el retorno acumulado descontado. La Q-function es el puente que conecta estado, acción y retorno, mientras que la política es la regla de conducta del Agent.


